\chapter{Introdução}
\label{chap:introducao}

A produção de mel e de outros produtos apícolas integra a atividade agropecuária brasileira \cite{ibge2024}. No contexto da pesquisa, a identificação da origem de uma amostra e sua associação aos resultados laboratoriais são necessárias para que os dados possam ser recuperados e interpretados posteriormente. A rastreabilidade pressupõe a possibilidade de acompanhar um produto e suas informações ao longo das etapas relevantes da cadeia \cite{iso22005}. Sistemas de informação laboratoriais oferecem uma forma de organizar esses registros \cite{gibbon1996}.

O NAPI Abelhas reúne pesquisadores de diferentes áreas e instituições no Paraná. O projeto apresentado no TCC 1 partiu da necessidade de organizar informações sobre produtores, locais de coleta, amostras e análises que estavam distribuídas em planilhas. A proposta definiu uma plataforma web e uma arquitetura inicial. O presente trabalho relata a implementação dessa plataforma, as escolhas efetivamente incorporadas aos projetos de software e as diferenças observadas em relação ao planejamento.

\section{Problema e objetivos}
\label{sec:objetivos}

O problema abordado é a fragmentação dos registros de amostras e resultados de análises, que dificulta associar cada resultado à sua origem e recuperar as informações para consultas posteriores. O objetivo geral deste trabalho é implementar e documentar uma plataforma web para o gerenciamento desses registros no contexto do NAPI Abelhas.

Os objetivos específicos são:
\begin{itemize}
    \item modelar produtores, pontos de coleta, espécies de abelhas, tipos de amostra, amostras, análises e arquivos associados;
    \item implementar uma interface de cadastro e consulta e uma API para persistência e recuperação dos dados;
    \item integrar autenticação e controle de acesso por organização e papel;
    \item disponibilizar o envio e a recuperação de arquivos associados às análises;
    \item comparar a implementação com os requisitos do TCC 1 e registrar as limitações constatadas.
\end{itemize}

\section{Procedimentos adotados}
\label{sec:metodologia}

O trabalho tem caráter aplicado: parte de uma proposta de sistema e apresenta seus artefatos de implementação. A descrição técnica foi construída por inspeção do código-fonte do frontend, da API, do esquema Prisma e da configuração de infraestrutura. A comparação com o TCC 1 considera uma funcionalidade implementada apenas quando há evidência correspondente nesses artefatos. A presença de uma tecnologia no arquivo de dependências ou na configuração de contêineres, isoladamente, não demonstra que a funcionalidade planejada esteja em operação.

A verificação descrita no Capítulo \ref{chap:resultados} separa a inspeção estática, os testes automatizados disponíveis e as capturas de interface previstas. Não houve estudo de usabilidade, entrevista ou validação com pesquisadores e técnicos do NAPI Abelhas. Consequentemente, o texto não atribui à plataforma ganhos mensurados de produtividade, desempenho ou qualidade dos dados.

\section{Organização do trabalho}

O Capítulo \ref{chap:sistema_implementado} apresenta requisitos, atividades de desenvolvimento, arquitetura e persistência, indicando as alterações de escopo observadas. O Capítulo \ref{chap:resultados} reúne as funcionalidades identificadas, as evidências técnicas e as limitações. O Capítulo \ref{chap:conclusao} sintetiza as contribuições e indica trabalhos futuros.
